\clearpage
\section{Biến đổi tuyến tính, phân phối biên và tính độc lập}
\label{sec:linear}

\subsection{Một công thức kỳ vọng, một công thức hiệp phương sai}
Cho vectơ ngẫu nhiên $\mathbf X$ có trung bình $\vmu$ và hiệp phương sai
$\Sigma$, đặt $\mathbf Y=A\mathbf X+\mathbf d$, trong đó $A$ và
$\mathbf d$ là các hằng số. Khi đó
\begin{equation}
 \E(\mathbf Y)=A\vmu+\mathbf d,\qquad
 \Cov(\mathbf Y)=A\Sigma A'.
 \label{eq:linearmeanvar}
\end{equation}
Hai công thức về moment này đúng ngay cả khi $\mathbf X$ không chuẩn, miễn
các moment tồn tại. Để kết luận \emph{phân phối của $\mathbf Y$ cũng chuẩn},
ta cần giả thiết chuẩn nhiều chiều của $\mathbf X$.

Tại sao $\mathbf d$ không xuất hiện trong hiệp phương sai? Vì
\[
 \mathbf Y-\E(\mathbf Y)=A(\mathbf X-\vmu).
\]
Tịnh tiến mọi giá trị cùng một lượng chỉ đổi tâm, không đổi độ phân tán.
Khi lấy tích ngoài và kỳ vọng, ta được
\[
 \E\bigl[A(\mathbf X-\vmu)(\mathbf X-\vmu)'A'\bigr]=A\Sigma A'.
\]

\subsection{Tổ hợp tuyến tính một biến}
\begin{result}
Nếu $\mathbf X\sim\N_p(\vmu,\Sigma)$, thì với $\mathbf a\in\R^p$,
\begin{equation}
 \mathbf a'\mathbf X\sim\N(\mathbf a'\vmu,\mathbf a'\Sigma\mathbf a).
 \label{eq:unilinear}
\end{equation}
Trường hợp $\mathbf a=0$ cho biến hằng, tức chuẩn suy biến với phương sai 0.
\end{result}
\nguon{Chương 4, §4.2, Result 4.2, tr.~156}

Viết công thức phương sai ra từng thành phần sẽ thấy điều dễ bị bỏ sót:
\[
 \Var\!\left(\sum_i a_iX_i\right)
 =\sum_i a_i^2\sigma_{ii}+2\sum_{i<k}a_i a_k\sigma_{ik}.
\]
Không được chỉ cộng các phương sai nếu các biến chưa được chứng minh không
tương quan. Các hệ số phải được bình phương ở các hạng đường chéo.

\begin{example}[Tổ hợp từ dữ kiện BTVN 1]
Dùng bộ tham số
\begin{equation}
 \vmu=\begin{pmatrix}1\\-1\\2\end{pmatrix},\qquad
 \Sigma=\begin{pmatrix}4&0&-1\\0&5&0\\-1&0&2\end{pmatrix}.
 \label{eq:running3d}
\end{equation}
Bộ này xuất hiện trong Bài 2.1--2.3 của BTVN 1. Ta tự xây dựng tổ hợp
$Y=3X_1-2X_2+X_3$ để minh họa; không dùng lẫn nó với bộ tham số khác của Bài 2.4.

\buoc{1}{Ghi vectơ hệ số}
$\mathbf a=(3,-2,1)'$.
\buoc{2}{Tính trung bình}
$\E(Y)=3(1)-2(-1)+2=7$.
\buoc{3}{Tính phương sai}
\[
 \Sigma\mathbf a=\begin{pmatrix}11\\-10\\-1\end{pmatrix},\qquad
 \mathbf a'\Sigma\mathbf a=3(11)+(-2)(-10)+1(-1)=52.
\]
Cách khai triển cho cùng kết quả:
$9\cdot4+4\cdot5+2+2\cdot3\cdot1\cdot(-1)=52$.
\buoc{4}{Kết luận cả dạng phân phối và tham số}
$Y\sim\N(7,52)$; độ lệch chuẩn là $\sqrt{52}$, không phải 52.
\end{example}

\subsection{Nhiều tổ hợp cùng lúc}
\begin{result}
Nếu $\mathbf X\sim\N_p(\vmu,\Sigma)$ và $A$ có kích thước $q\times p$,
thì
\begin{equation}
 A\mathbf X+\mathbf d\sim\N_q(A\vmu+\mathbf d,A\Sigma A').
 \label{eq:multilinear}
\end{equation}
Nếu $\Sigma$ xác định dương và $A$ có hạng hàng đầy đủ $q$, hiệp phương
sai mới xác định dương. Nếu $A$ thiếu hạng, kết quả là chuẩn suy biến,
không có mật độ thông thường trên toàn bộ $\R^q$.
\end{result}
\nguon{Chương 4, §4.2, Result 4.3, tr.~157; điều kiện hạng được giải thích thêm ở đây}

\begin{example}[Tổng và hiệu]
Với $\mathbf X\sim\N_2((10,20)',\begin{pmatrix}4&3\\3&9\end{pmatrix})$,
đặt $Y_1=X_1-X_2$, $Y_2=X_1+X_2$. Khi đó
\[
 A=\begin{pmatrix}1&-1\\1&1\end{pmatrix},\qquad
 A\vmu=\begin{pmatrix}-10\\30\end{pmatrix}.
\]
Ta tính
\[
 \Var(Y_1)=4+9-2(3)=7,\quad
 \Var(Y_2)=4+9+2(3)=19,
\]
\[
 \Cov(Y_1,Y_2)=\Var(X_1)-\Var(X_2)=4-9=-5.
\]
Vậy $(Y_1,Y_2)'\sim\N_2((-10,30)',\begin{pmatrix}7&-5\\-5&19\end{pmatrix})$.
Tổng và hiệu chưa độc lập; tương quan giữa các biến gốc không phải yếu tố
duy nhất quyết định hiệp phương sai giữa tổng và hiệu.
\end{example}

\subsection{Phân phối biên: chọn biến, không điều kiện hóa}
Phân phối biên mô tả các biến được giữ lại \emph{khi không biết giá trị}
của các biến bị bỏ. Đối với chuẩn nhiều chiều, chỉ cần chọn đúng phần tử
của $\vmu$ và $\Sigma$.

\begin{result}
Nếu phân hoạch
\[
 \mathbf X=\begin{pmatrix}\mathbf X_1\\\mathbf X_2\end{pmatrix},\quad
 \vmu=\begin{pmatrix}\vmu_1\\\vmu_2\end{pmatrix},\quad
 \Sigma=\begin{pmatrix}\Sigma_{11}&\Sigma_{12}\\\Sigma_{21}&\Sigma_{22}\end{pmatrix},
\]
thì $\mathbf X_1\sim\N_{q_1}(\vmu_1,\Sigma_{11})$ và tương tự cho
$\mathbf X_2$.
\end{result}
\nguon{Chương 4, §4.2, Result 4.4, tr.~158--159}

\begin{example}[Chọn các thành phần không liền nhau]
Trong \eqref{eq:running3d}, phân phối của $(X_1,X_3)'$ là
\[
 \N_2\!\left(\begin{pmatrix}1\\2\end{pmatrix},
 \begin{pmatrix}4&-1\\-1&2\end{pmatrix}\right).
\]
Ta chọn hàng và cột 1, 3 theo \emph{cùng thứ tự}. Nếu chọn $(X_3,X_1)'$,
trung bình phải là $(2,1)'$ và hiệp phương sai là
$\begin{pmatrix}2&-1\\-1&4\end{pmatrix}$.
\end{example}

Không cần tích phân mật độ để giải dạng bài này. Một ma trận chọn
$A=\begin{pmatrix}1&0&0\\0&0&1\end{pmatrix}$ trong
\eqref{eq:multilinear} cho ngay kết quả.

\subsection{Độc lập: vai trò quyết định của chuẩn đồng thời}
\begin{result}
Với hai khối thuộc một vectơ \emph{chuẩn đồng thời},
\begin{equation}
 \mathbf X_1\perp\mathbf X_2
 \quad\Longleftrightarrow\quad \Sigma_{12}=0.
 \label{eq:normalindependence}
\end{equation}
Với hai tổ hợp $U=\mathbf a'\mathbf X$, $V=\mathbf b'\mathbf X$,
\begin{equation}
 U\perp V\quad\Longleftrightarrow\quad
 \mathbf a'\Sigma\mathbf b=0.
 \label{eq:linearindependence}
\end{equation}
\end{result}
\nguon{Chương 4, §4.2, Result 4.5, tr.~159--160; phát biểu cho tổ hợp là hệ quả của Results 4.3 và 4.5}

\hieu{Trong mô hình chuẩn, khi khối hiệp phương sai chéo bằng không,
mật độ chung tách được thành tích hai mật độ biên. Đó chính là tính độc lập.
Ngoài mô hình chuẩn, thông tin hiệp phương sai không đủ để mô tả mọi kiểu phụ thuộc.}

Trong \eqref{eq:running3d}, $\Cov(X_1,X_2)=0$ nên $X_1\perp X_2$;
$\Cov(X_1,X_3)=-1$ nên $X_1$ và $X_3$ không độc lập. Hơn nữa,
\[
 \Cov\bigl(X_2,(X_1,X_3)'\bigr)=(0,0)
\]
nên $X_2$ độc lập với \emph{cả vectơ} $(X_1,X_3)'$.
Trong chuẩn nhiều chiều, kết luận này được bảo đảm bởi khối hiệp phương sai.
Trong phân phối tổng quát, độc lập từng cặp không nhất thiết đủ để suy ra
độc lập với cả vectơ.

\begin{example}[Không tương quan nhưng vẫn phụ thuộc]
Cho $Z\sim\N(0,1)$ và $W=Z^2$. Ta có
$\Cov(Z,W)=\E(Z^3)-\E(Z)\E(Z^2)=0$.
Nhưng $W$ do $Z$ quyết định nên không độc lập với $Z$.
Cặp $(Z,W)$ không chuẩn nhiều chiều.
Ví dụ này cho thấy không thể áp dụng \eqref{eq:normalindependence}
chỉ vì một biến có phân phối chuẩn.
\end{example}

Ngay cả việc mỗi biến riêng lẻ đều chuẩn cũng chưa đủ. Cho $Z\sim\N(0,1)$,
cho $B$ độc lập với $Z$, nhận $\pm1$ với xác suất $1/2$; đặt $V=BZ$.
Cả $Z,V$ đều chuẩn và $\Cov(Z,V)=0$, nhưng $|V|=|Z|$ nên chúng phụ thuộc.
Chuẩn \emph{đồng thời} là giả thiết mạnh hơn chuẩn của từng biên.

\subsection{Tìm hệ số để triệt tiêu hiệp phương sai}
\begin{example}[Bài 2.4b của BTVN 1]
Ở bài này, bộ tham số khác với \eqref{eq:running3d}:
\[
 \vmu=(2,-3,1)',\qquad
 \Sigma=\begin{pmatrix}1&1&1\\1&3&2\\1&2&2\end{pmatrix}.
\]
Đặt $\mathbf Z=(X_1,X_3)'$ và $W=X_2-a_1X_1-a_2X_3$.
Đề yêu cầu $X_2\perp W$. Vì cặp này chuẩn đồng thời, ta giải
\[
 \Cov(X_2,W)=3-a_1-2a_2=0.
\]
Toàn bộ nghiệm là $a_1=3-2t$, $a_2=t$, $t\in\R$.
Đây là một phương trình với hai ẩn, nên không có nghiệm duy nhất.
\end{example}

\luuy{Yêu cầu $X_2\perp W$ khác với yêu cầu $\mathbf Z\perp W$.
Trong hồi quy ở phần kế tiếp, ta chọn hệ số để phần dư độc lập với
\emph{các biến dùng dự báo}, chứ không phải với biến cần dự báo.
Hãy đọc đúng hai đại lượng mà đề yêu cầu độc lập.}

\kiemtra{Với chuẩn hai chiều, $\Var(X_1)=\Var(X_2)$, tổng và hiệu có độc
lập không? Có: hiệp phương sai của chúng bằng hiệu hai phương sai, tức 0.}
